\uuid{E4Pa}
\exo7id{3826}
\titre{exo7 3826}
\auteur{quercia}
\organisation{exo7}
\datecreate{2010-03-11}
\isIndication{false}
\isCorrection{false}
\chapitre{Espace euclidien, espace normé}
\sousChapitre{Espaces vectoriels hermitiens}
\cpge{mpsi-mp2i.applications-lineaires.e, pcsi.applications-lineaires.f, mpsi-mp2i.espaces-prehilbertiens.a, pcsi.espaces-prehilbertiens.a, ptsi.geometrie-plan.b, mp-mpi.topologie-espaces-normes.a, pc.espaces-vectoriels-normes.a, psi.espaces-vectoriels-normes.a, psi.espaces-prehilbertiens.a, pt.structure-prehilbertienne.a}
\module{Algèbre}
\niveau{L2}
\difficulte{}

\contenu{
\question{
Soit $E$ un espace préhilbertien réel et $B$ une boule ouverte de $E$ ne
contenant pas $\vec 0$.
Montrer qu'il existe une forme linéaire $f \in E^*$ telle que :
$\forall\ \vec x \in B,\ f(\vec x) > 0$.
}
}
